\documentclass[10pt,a4paper]{amsart}
\usepackage[margin=19mm]{geometry}
\usepackage{amsmath,amssymb,mathtools}
\usepackage[scr=rsfs,cal=euler]{mathalfa}
\usepackage{xfrac}
\usepackage[unicode,hidelinks,bookmarks=false]{hyperref}
\newcommand{\ie}{i.e.\ }
\newcommand{\cf}{cf.\ }
\newcommand{\resp}{respectively\ }
\newcommand{\Subsection}{\subsection}
\providecommand{\nobreakdash}{\nobreak\hbox{-}}
\setlength{\emergencystretch}{2em}
\multlinegap0pt
\usepackage{tikz}
\usepackage{amssymb}
\newtheorem{theorem}{Theorem}
\newtheorem{lemma}{Lemma}
\title{Farey graphs and geodesic expansions of complex continued fractions}
\author{Hitoshi Nakada, Rie Natsui,, J\"org Thuswaldner}
\date{Source: arXiv:2603.28468v2 (CC BY 4.0)}
\begin{document}
\maketitle

\noindent\emph{Source and licence.} Farey graphs and geodesic expansions of complex continued fractions. Version arXiv:2603.28468v2 (CC BY 4.0), distributed by arXiv under the Creative Commons Attribution 4.0 International licence, http://creativecommons.org/licenses/by/4.0/, as recorded in the arXiv licence metadata for this exact version and shown on the abstract page https://arxiv.org/abs/2603.28468v2. Full licence notice: \texttt{licenses/arxiv-2603.28468-CC-BY-4.0.txt}.\par\smallskip

\noindent\textit{Editorial scope.} Counted unit: the article's theorem theorem (source0.tex lines 950--952). The article's complete original proof of it is delivered with it (source0.tex lines 954--1008, one \texttt{proof} environment of 55 lines). No mathematics is written, repaired or re-derived: the statement and the proof are the article's own lines, in the article's own notation, and no retained line is substituted or re-flowed.  Retained uncounted as the source-local context the proof needs: lines 307--310; lines 936--940.  Nothing was omitted from any retained span, so the package carries no omission record.  No retained line contains a citation, so the wrapper carries no bibliography and no bibliography entry.  No retained line contains a figure environment, an image-inclusion command or drawing code, so no image is delivered and the package declares no asset dependency.  Editorial formatting only, in addition: an AMS \texttt{amsart} preamble that restates the article's own declarations and macros used by the retained lines, namely the environments \texttt{theorem}, \texttt{lemma} on separate counters, together with the standard packages those lines need.  The retained environments print in this document as Theorem 1, Lemma 1 in order of appearance, whereas the article prints the same items with the numbering of its own full document.  The complete native source of the article as distributed in the arXiv e-print for this exact version is bundled unchanged as \texttt{sources/197444/source0.tex}.
\begin{theorem}\label{th:correspondencePathCF}
Fix $d\in\{1,2,3,7,11\}$ and let $x\in \mathbb{Q}(\sqrt{-d})$. Then $R_1,\ldots, R_n$ with $R_n=x$, are the consecutive convergents of some complex continued fraction expansion of $x$ if and only if $\infty \to  R_1 \to \cdots \to  R_n$ is a walk in $\mathcal{E}_{d}$ from 
$\infty$ to $x$.
\end{theorem}

\begin{lemma}\label{lem:prepgeoalg}
For $d\in \{1,2,3\}$ let $z \in \mathbb Q(\sqrt{-d})$.  Then the reflection $\hat{z}$ of $z$ across 
$L_{d}$ is contained in $\mathbb Q(\sqrt{-d})$ and the Ford spheres $\mathcal F_d(z)$ 
and $\mathcal F_d(\hat{z})$ have the same radius.  
\end{lemma} 

\begin{theorem}
Let $d \in \{1, 2, 3\}$. Then the nearest integer complex continued fraction map~$T_d$ always produces geodesic continued fraction expansions. 
\end{theorem}

\begin{proof}
Let $\frac{P}{Q} \in U_{d}$ be given. Suppose that
$\frac{P}{Q}= [a_1,\ldots, a_n]$ is the nearest integer continued fraction expansion of 
$\frac{P}{Q}$. If $\frac{P_k}{Q_k}$ is the $k$-th convergent of $\frac PQ$, in view of Theorem~\ref{th:correspondencePathCF} this expansion corresponds to the walk
\begin{equation}\label{eq:HurwitzPath}
\infty=\frac{P_{-1}}{Q_{-1}} \to 0=\frac{P_0}{Q_0} \to \frac{P_1}{Q_1} \to \frac{P_2}{Q_2} \to \cdots \to \frac{P_n}{Q_n}=\frac{P}{Q}
\end{equation}
in the Farey graph $\mathcal{E}_{d}$. Now suppose that
$\infty \to r_0 \to r_1 \to \cdots \to r_{\ell}=\frac{P}{Q}$ is a geodesic path in $\mathcal{E}_{d}$. We need to show that \eqref{eq:HurwitzPath} is also a geodesic path in $\mathcal{E}_{d}$. If these two walks are equal there is nothing to do. If not, we choose $m$ such that $\frac{P_{k}}{Q_{k}}=r_{k}$ holds for $0 \leq k <m$, and $\frac{P_m}{Q_m} \neq r_m$. Thus $\infty \to \frac{P_0}{Q_0}\to\cdots\to \frac{P_{m-1}}{Q_{m-1}}$ is geodesic. Consider the linear fractional map 
$$
S=\begin{pmatrix}
0 & 1 \\
1 & a_{m-1}
\end{pmatrix}^{-1}
\begin{pmatrix}
0 & 1 \\
1 & a_{m-2}
\end{pmatrix}^{-1} 
\cdots 
\begin{pmatrix}
0 & 1 \\
1 & a_1
\end{pmatrix}^{-1}
$$
($S$ is the identity if $m=0$).  Then $S\big(\frac{P_{m-2}}{Q_{m-2}}\big)=S(r_{m-2})=\infty$, $S\big(\frac{P_{m-1}}{Q_{m-1}}\big)=S(r_{m-1})=0$, and $S\big(\frac{P_m}{Q_m}\big)=\frac{1}{a_m}$. 
Since M\"obius transformations preserve adjacency in $\mathcal{E}_d$ this implies that $0 \to S\left(r_m\right)$ and, hence, $\infty \to S\left(r_m\right)^{-1}$. Thus
$S\left(r_m\right)^{-1}\in\mathbb Z[\omega_{d}]$. By the choice of $m$ we have  $S\left(r_m\right)^{-1}\not=a_m$ and, hence, $S\left(r_m\right)^{-1}-a_m \notin U_d$. 
Writing $r_k^{(1)}=S\left(r_k\right)^{-1}-a_m$ for $k \geq m$ we see that  
\begin{equation}\label{eq:rprimepath}
\infty \to r_m^{(1)} \to \cdots \to r_\ell^{(1)}
\end{equation}
is a geodesic path in $\mathcal{E}_d$. We have $r_m^{(1)}\not\in{U_{d}}^{\mathrm{cl}}$ and  
\[
r_{\ell}^{(1)}=S(r_\ell)^{-1}-a_m= S\big(\frac PQ\big)^{-1}-a_m= T_{d}^{m-1}\big(\frac PQ\big)^{-1}-a_m = T_{d}^m\big(\frac PQ\big) \in U_{d}.
\]
Thus there exists $t\in\{m+1,\ldots,\ell\}$ such that $r_s^{(1)}\not\in {U_{d}}^{\mathrm{cl}}$ for $m \leq s<t$ and $r_t^{(1)}\in {U_{d}}^{\mathrm{cl}}$.  The following three cases can occur.
\begin{itemize}
\item[(a)]\,  $r_t^{(1)}$ is symmetric to $r_{t-1}^{(1)}$ w.r.t. some line $L_d$ containing a line segment of $\partial U_{d}$.
\item[(b)] \, $r_t^{(1)}$ not on $\partial U_{d}$  and not symmetric to $r_{t-1}^{(1)}$ w.r.t.\ any line $L_d$ containing a line segment of $\partial U_{d}$.
\item[(c)]\, $r_t^{(1)}$ is on some line $L_d$ containing a line segment of $\partial U_{d}$.
\end{itemize}

Case (a): For each $s\in\{m,\ldots, t-1\}$ let $r_s^{(2)}$ be symmetric to $r_s^{(1)}$ w.r.t.\ $L_d$. Then $r_{t-1}^{(2)}=r_t^{(1)}$. By Lemma~\ref{lem:prepgeoalg}, this yields the walk 
$\infty\to r_m^{(2)}\to \cdots \to r_{t-1}^{(2)}=r_t^{(1)}\to r_{t+1}^{(1)} \to \cdots \to r_\ell^{(1)}$ that is shorter than the geodesic path in \eqref{eq:rprimepath}, a contradiction.

Case (b): We consider the projections of the Ford spheres $G_{t-1}$ and $G_t$ associated with $r_{t-1}^{(1)}$ and $r_t^{(1)}$ to $\mathbb{C}$. These projections are circles. These circles cannot cross a line $L_d\subset \partial U_{d}$. Indeed, if the Ford sphere $G$ of some point $p\not\in L_d$ would cross $L_d$, it would intersect the Ford sphere $G'$ of the reflection of $p$ across $L_d$ (because, by Lemma~\ref{lem:prepgeoalg}, $G$ and $G'$ have the same radius), a contradiction.  So the Ford spheres $G_{t-1}$ and $G_t$ cannot be tangent to each other. Thus $r_{t-1}^{(1)}$ and $r_t^{(1)}$ are not adjacent in $\mathcal{E}_d$, a contradiction. 

Thus only case (c) can occur. In this case by Lemma~\ref{lem:prepgeoalg} we can construct a walk $\infty\to r_1^{(2)}\to \cdots\to r_t^{(2)}=r_t^{(1)}$, where $r_s^{(2)}$ is symmetric to $r_s^{(1)}$ w.r.t.\ $L_d$ for $s\in \{1,\ldots, t\}$. This walk has the same length as  the geodesic path $\infty\to r_1^{(1)}\to \cdots\to r_t^{(1)}$. Setting $r_k^{(2)}=r_k^{(1)}$ for $t \le k \le \ell$ yields a geodesic path $\infty \to r_m^{(2)}\to \cdots \to r_\ell^{(2)}$ and $t_2<t$ such that $r_s^{(2)} \not\in {U_{d}}^{\mathrm{cl}}$ for $m \leq s<t_2$ and $r_{t_2}^{(2)}\in {U_{d}}^{\mathrm{cl}}$. Iterating this procedure yields a geodesic path $\infty \to r_m^{(n)} \to \cdots \to r_\ell^{(n)}$ such that  $r_{m}^{(n)}\in{U_d}^{\mathrm{cl}}$, {\em i.e.} $r_m^{(n)}=0$. Thus
\[
\begin{split}
&\infty \to r_0 \to r_1 \to \cdots \to r_{m-1} \to S^{-1}\bigg(\frac1{a_m + r_m^{(n)}}\bigg)  \to  \cdots \to S^{-1}\bigg(\frac1{a_m + r_\ell^{(n)}}\bigg)
\end{split}
\]
is a  geodesic path from $\infty$ to $S^{-1}\bigg(\frac1{a_m + r_\ell^{(n)}}\bigg)=r_\ell=\frac PQ$ with $r_0=\frac {P_0}{Q_0},\ldots, r_{m-1}=\frac {P_{m-1}}{Q_{m-1}}$ and $S^{-1}\bigg(\frac1{a_m + r_m^{(n)}}\bigg)=\frac{P_m}{Q_m}$. Thus, in particular, $\infty \to \frac{P_0}{Q_0}\to\cdots\to \frac{P_{m}}{Q_{m}}$ is geodesic. The fact that \eqref{eq:HurwitzPath} is geodesic now follows by induction on $m$.
\end{proof}

\end{document}
